\subsection{Independent Variable}

The single independent variable across the two experimental conditions is
the presence or absence of retrieved evidence in the prompt (baseline vs.\
retrieval-augmented; Section~\ref{sec:method}). All other factors --- model
checkpoint, decoding parameters, test examples, and scoring code --- are
held fixed between conditions by construction.

\subsection{Model and Decoding Configuration}

The baseline checkpoint is \texttt{Salesforce/blip2-opt-2.7b}, with the
vision encoder frozen and the language model unfrozen
(\texttt{configs/model\_config.yaml}); \texttt{microsoft/llava-med-v1.5-mistral-7b}
is retained as an alternate candidate to substitute if the primary
checkpoint underperforms during initial validation. Generation is
deterministic (greedy decoding, \texttt{do\_sample=false}) with a maximum of
128 new tokens, so that repeated runs of either condition on the same input
are reproducible. The retriever uses \texttt{openai/clip-vit-base-patch32}
as its image encoder and returns the top $k{=}5$ most similar corpus items
per query by default; Section~\ref{sec:results} additionally reports a
sweep over $k \in \{1, 3, 5, 10\}$ to test the sensitivity of accuracy to
this choice.

\subsection{Evaluation Metrics}

VQA-RAD questions are scored differently by type, following the
distinction the dataset itself makes between closed-ended and open-ended
questions~\cite{lau2018vqarad}: closed-ended (yes/no-style) questions are
scored by exact match after normalization (lowercasing, article and
punctuation removal); open-ended free-text questions are scored by BLEU-4
(with smoothing, appropriate given the short length of typical VQA-RAD
answers) and ROUGE-L F-measure. Both conditions are scored by the same
metric implementation, so the two accuracy numbers in
Section~\ref{sec:results} are computed identically.

\subsection{Statistical Significance}

Given VQA-RAD's test split is small (on the order of a few hundred
examples), a raw accuracy delta between conditions could plausibly reflect
sampling noise rather than a genuine effect. This is tested with a paired
bootstrap: per-example scores from both conditions are resampled with
replacement 10{,}000 times, preserving the pairing between a given test
example's baseline and RAG scores, and a two-sided p-value together with a
95\% confidence interval is reported for the mean delta. An effect is
treated as reliable only when $p < 0.05$ and the confidence interval
excludes zero.

\subsection{Qualitative Error Analysis}

Beyond the aggregate delta, every shared test example is categorized as
\emph{improved} (RAG correct, baseline incorrect), \emph{regressed}
(baseline correct, RAG incorrect), or unchanged (both correct or both
incorrect), to characterize where retrieval helps or hurts rather than only
whether it does so on average. Retrieved-evidence quality is profiled
separately from answer correctness --- mean top-1 retrieval similarity and
the proportion of queries for which no evidence was retrieved at all --- so
that a weak result can be attributed to the retriever failing to find
relevant images versus the model failing to make use of evidence it was
given.

\subsection{Hardware and Reproducibility}

All experiments were run on a Kaggle Notebook with a single NVIDIA T4 GPU
(15GB VRAM), using PyTorch 2.3.1 (with \texttt{transformers} upgraded to
4.46.3 and \texttt{tokenizers} to $\geq$0.20.3 to resolve a tokenizer
deserialization incompatibility with the pinned \texttt{transformers}
version) and faiss-cpu for index construction and search. The baseline
vision-language model was loaded in half precision (float16) to fit
comfortably within available GPU memory. Corpus embedding and index
construction over the full ROCOv2 corpus (59{,}962 image-caption pairs
after filtering) took approximately 17 minutes; each full evaluation run
over the 451-example VQA-RAD official test split (baseline or a single
retrieval-augmented condition) took approximately 2--3 minutes. Every
number reported in Section~\ref{sec:results} is generated by the released
code (\texttt{scripts/run\_baseline.py}, \texttt{scripts/run\_rag\_eval.py},
\texttt{scripts/run\_experiments.py},
\texttt{scripts/generate\_paper\_results.py}) against a fixed, versioned
configuration (\texttt{configs/model\_config.yaml}), so any of these
figures can be regenerated by re-running the corresponding script.
